\documentclass[a4paper]{article}

% \usepackage[default]{fontsetup}

\usepackage{fancyhdr}
\usepackage{extramarks}
\usepackage{amsmath}
\usepackage{amsthm}
\usepackage{amsfonts}
\usepackage{tikz}
\usepackage[plain]{algorithm}
\usepackage{algpseudocode}
\usepackage{enumerate}
\usepackage{tikz}

\usetikzlibrary{automata,positioning}

%
% Basic Document Settings
%  

\topmargin=-0.2in
\evensidemargin=0in
\oddsidemargin=0in
\textwidth=6.5in
\textheight=9.5in
\headsep=0.25in

\linespread{1.1}

\pagestyle{fancy}
\lhead{\hmwkAuthorName}
\chead{\hmwkClass : \hmwkTitle}
\rhead{\firstxmark}
\lfoot{\lastxmark}
\cfoot{\thepage}

\renewcommand\headrulewidth{0.4pt}
\renewcommand\footrulewidth{0.4pt}

\setlength\parindent{0pt}

%
% Create Problem Sections
%

\newcommand{\enterProblemHeader}[1]{
    \nobreak\extramarks{}{Problem \arabic{#1} continued on next page\ldots}\nobreak{}
    \nobreak\extramarks{Problem \arabic{#1} (continued)}{Problem \arabic{#1} continued on next page\ldots}\nobreak{}
}

\newcommand{\exitProblemHeader}[1]{
    \nobreak\extramarks{Problem \arabic{#1} (continued)}{Problem \arabic{#1} continued on next page\ldots}\nobreak{}
    \stepcounter{#1}
    \nobreak\extramarks{Problem \arabic{#1}}{}\nobreak{}
}

\newcommand*\circled[1]{\tikz[baseline=(char.base)]{
		\node[shape=circle,draw,inner sep=2pt] (char) {#1};}}


\setcounter{secnumdepth}{0}
\newcounter{partCounter}
\newcounter{homeworkProblemCounter}
\setcounter{homeworkProblemCounter}{1}
\nobreak\extramarks{Problem \arabic{homeworkProblemCounter}}{}\nobreak{}

%
% Homework Problem Environment
%
% This environment takes an optional argument. When given, it will adjust the
% problem counter. This is useful for when the problems given for your
% assignment aren't sequential. See the last 3 problems of this template for an
% example.
%

\newenvironment{homeworkProblem}[2][-1]{%
    \ifnum#1>0\relax
        \setcounter{homeworkProblemCounter}{#1}
    \fi
    \section{Problem \arabic{homeworkProblemCounter} (#2 points)}
    \setcounter{partCounter}{1}
    \enterProblemHeader{homeworkProblemCounter}
}{%
    \exitProblemHeader{homeworkProblemCounter}
}

%
% Homework Details
%   - Title
%   - Class
%   - Due date
%   - Name
%   - Student ID

\newcommand{\hmwkTitle}{Homework\ \#08}
\newcommand{\hmwkClass}{Probability \& Statistics for EECS}
\newcommand{\hmwkDueDate}{2026-11-22}
\newcommand{\hmwkAuthorName}{Zhang Yiheng}
\newcommand{\hmwkAuthorID}{2025531068}


%
% Title Page
%

\title{
    \vspace{2in}
    \textmd{\textbf{\hmwkClass:\\  \hmwkTitle}}\\
    \normalsize\vspace{0.1in}\small{Due\ on\ \hmwkDueDate\ at 23:59}\\
	\vspace{4in}
}

\author{
	Name: \textbf{\hmwkAuthorName} \\
	Student ID: \hmwkAuthorID}
\date{}

\renewcommand{\part}[1]{\textbf{\large Part \Alph{partCounter}}\stepcounter{partCounter}\\}

%
% Various Helper Commands
%

% Useful for algorithms
\newcommand{\alg}[1]{\textsc{\bfseries \footnotesize #1}}
% For derivatives
\newcommand{\deriv}[1]{\frac{\mathrm{d}}{\mathrm{d}x} (#1)}
% For partial derivatives
\newcommand{\pderiv}[2]{\frac{\partial}{\partial #1} (#2)}
% Integral dx
\newcommand{\dx}{\mathrm{d}x}
% Alias for the Solution section header
\newcommand{\solution}{\textbf{\large Solution}}
% Probability commands: Expectation, Variance, Covariance, Bias
\newcommand{\E}{\mathrm{E}}
\newcommand{\Var}{\mathrm{Var}}
\newcommand{\Cov}{\mathrm{Cov}}
\newcommand{\Bias}{\mathrm{Bias}}

\begin{document}


% \maketitle
% \thispagestyle{empty}
% \pagebreak

\date{
Due on Nov. 22, 2026, 23:59 UTC+8}
\title{SI 140A.02  Probability \& Statistics for EECS, Fall 2026 \\
Homework 8}
\author{Name: \textbf{\hmwkAuthorName} \\ Student ID: \hmwkAuthorID}
\maketitle
Read all the instructions below carefully before you start working on the assignment, and before you make a submission.
\begin{itemize}
    \item Submit your homework as a single PDF file.
Written solutions to Problems 1, 2, 4, 5, and 6 must be
prepared using \textbf{LaTeX}. For Problem 3, complete the
provided Jupyter notebook and export it to PDF.
Submissions that do not follow these format requirements
will receive no credit.
    \item You are required to write down all the major steps towards making your conclusions; otherwise you may obtain limited points of the problem.
    \item Write your homework in English; otherwise you will get no points of this homework.
    \item Any form of plagiarism will lead to $0$ point of this homework. 
\end{itemize}
\newpage

\begin{homeworkProblem}{10}

Let $X_1,X_2,X_3,X_4$ be independent and identically distributed
continuous random variables. Their distribution is unspecified.
Prove that
\[
\Pr(X_1<X_2,\ X_3<X_2,\ X_3<X_4)=\frac{5}{24}.
\]

\medskip
\solution


Since $X_1,\dots,X_4$ are i.i.d.\ and continuous, $\Pr(X_i=X_j)=0$ for every
$i\neq j$, so with probability $1$ the four values are distinct and determine a
strict ordering. Because the joint distribution of $(X_1,X_2,X_3,X_4)$ is
invariant under any permutation of the indices (i.i.d.\ assumption), each of the
$4!=24$ possible orderings has the same probability $\frac1{24}$, regardless of
the common distribution.

The event $A=\{X_1<X_2,\ X_3<X_2,\ X_3<X_4\}$ is a union of orderings, so
$\Pr(A)=\frac{\#\{\text{orderings satisfying }A\}}{24}$. We count orderings
(from smallest to largest) consistent with $1<2$, $3<2$, $3<4$. The smallest
element must have nothing required below it, so it is $X_1$ or $X_3$.
\begin{itemize}
\item Smallest is $X_1$: the remaining $X_2,X_3,X_4$ must satisfy $X_3<X_2$ and
$X_3<X_4$, so $X_3$ is next and $X_2,X_4$ can be in either order:
$2$ orderings ($1\,3\,2\,4$ and $1\,3\,4\,2$).
\item Smallest is $X_3$: the remaining $X_1,X_2,X_4$ only need $X_1<X_2$, which
holds in exactly half of the $3!=6$ orderings: $3$ orderings
($3\,1\,2\,4$, $3\,1\,4\,2$, $3\,4\,1\,2$).
\end{itemize}
Hence there are $2+3=5$ favourable orderings and
\[
\Pr(X_1<X_2,\ X_3<X_2,\ X_3<X_4)=\frac{5}{24}.
\]

\end{homeworkProblem}
\newpage
\begin{homeworkProblem}{15}
Let $U_1,\ldots,U_n$ be independent $\operatorname{Unif}(0,1)$
random variables, where $n\ge1$. Define
\[
D=\min_{1\le i\le n}\left|U_i-\frac14\right|.
\]

\begin{enumerate}
\item[(a)] Find $\Pr(D>t)$ for every real $t$.
\item[(b)] Find $\mathbb E[D]$ without leaving any unevaluated integrals.
\end{enumerate}

\medskip
\solution


\textbf{(a)} Since $D\ge 0$, $\Pr(D>t)=1$ for $t<0$. For $t\ge0$, by independence,
\[
\Pr(D>t)=\Pr\Big(\bigcap_{i=1}^n\{|U_i-\tfrac14|>t\}\Big)=\Big(1-\Pr(|U_1-\tfrac14|\le t)\Big)^n .
\]
$\Pr(|U_1-\frac14|\le t)$ is the length of $[\frac14-t,\frac14+t]\cap[0,1]$:
\begin{itemize}
\item $0\le t<\frac14$: the interval lies inside $[0,1]$, length $2t$;
\item $\frac14\le t<\frac34$: it is $[0,\frac14+t]$, length $t+\frac14$;
\item $t\ge\frac34$: it covers $[0,1]$, length $1$.
\end{itemize}
Therefore
\[
\Pr(D>t)=
\begin{cases}
1, & t<0,\\[2pt]
(1-2t)^n, & 0\le t<\frac14,\\[2pt]
\left(\frac34-t\right)^n, & \frac14\le t<\frac34,\\[2pt]
0, & t\ge\frac34.
\end{cases}
\]
(The pieces agree at $t=\frac14$, both giving $2^{-n}$, and at $t=\frac34$, both giving $0$.)

\textbf{(b)} Since $D\ge0$, $\E[D]=\int_0^\infty\Pr(D>t)\,dt$. Thus
\[
\E[D]=\int_0^{1/4}(1-2t)^n\,dt+\int_{1/4}^{3/4}\left(\tfrac34-t\right)^n\,dt.
\]
With $s=1-2t$ in the first integral and $r=\frac34-t$ in the second,
\[
\int_0^{1/4}(1-2t)^n\,dt=\frac12\int_{1/2}^{1}s^n\,ds=\frac{1-2^{-(n+1)}}{2(n+1)},
\qquad
\int_{1/4}^{3/4}\left(\tfrac34-t\right)^n dt=\int_0^{1/2}r^n\,dr=\frac{2^{-(n+1)}}{n+1}.
\]
Adding,
\[
\E[D]=\frac{1-2^{-(n+1)}+2\cdot2^{-(n+1)}}{2(n+1)}
=\frac{1+2^{-(n+1)}}{2(n+1)}
=\frac{2^{n+1}+1}{(n+1)\,2^{n+2}}.
\]
Check ($n=1$): $\E|U-\frac14|=\frac12(\frac14)^2+\frac12(\frac34)^2=\frac{5}{16}$, and the
formula gives $\frac{5}{2\cdot 8}=\frac5{16}$.

\end{homeworkProblem}
\newpage
\begin{homeworkProblem}{15}
Complete Problem 3 in the accompanying Jupyter notebook.
The full problem statement, instructor-provided code, and
spaces for your written answers and implementation are
included in the notebook.

Complete all required written and coding sections, including
the derivations, algorithm descriptions, plots, numerical
comparisons, and explanations.

Export your completed notebook to PDF and include it as the
solution to Problem 3 in your homework PDF.

\medskip
\noindent\textbf{Converting the notebook to PDF.}
If you are unfamiliar with notebook-to-PDF conversion,
a Python script for converting an \texttt{.ipynb} file
to HTML will be provided. If you use PyCharm, you may follow these steps:
\begin{enumerate}
    \item Run all required code cells in your completed notebook
and save the notebook. Make sure that all required outputs,
plots, and written answers are visible before conversion.
    \item Run the provided Python script to convert your
    completed notebook to an HTML file.You need to open \texttt{export\_notebook.py} and replace
\texttt{your\_notebook.ipynb} with the path to your completed
notebook. Run the script using a Python environment with
\texttt{nbconvert} installed.
If needed, install it by running
\texttt{python -m pip install nbconvert} in the terminal.
    \item Right-click the generated HTML file and select
    \textbf{Open with}, then choose a web browser.
    \item In the browser, press \texttt{Ctrl+P}
    (\texttt{Command+P} on macOS), select
    \textbf{Save as PDF}, and save the file.
    \item Check that the saved PDF includes all required
    code, outputs, plots, and written answers, with no
    content cut off.
\end{enumerate}

\noindent\textbf{Submission.}
Use an online PDF merging tool 
like ilovepdf : https://www.ilovepdf.com to combine the PDF of your
written solutions and the PDF exported from your notebook
into \textbf{one PDF file}. Arrange the solutions in problem
order.

\medskip
\solution


See the completed notebook (exported PDF) attached after this page.

\end{homeworkProblem}
\newpage
\begin{homeworkProblem}{20}
Task A starts at time $0$ and requires
$X\sim\operatorname{Expo}(\lambda)$ units of time to finish.
Task B starts at a fixed time $a>0$ and requires
$Y\sim\operatorname{Expo}(\mu)$ units of time after its own start
to finish. Assume that $X$ and $Y$ are independent and
$\lambda,\mu>0$. Let
\[
T=\min\{X,a+Y\}
\]
be the time at which the first task finishes.

\begin{enumerate}
\item[(a)] Find $\mathbb E[T]$.
\item[(b)] Find the probability that task A finishes before task B.
\end{enumerate}

\medskip
\solution


\textbf{(a)} $T\ge0$, so $\E[T]=\int_0^\infty\Pr(T>t)\,dt$. By independence,
$\Pr(T>t)=\Pr(X>t)\Pr(a+Y>t)$, where $\Pr(X>t)=e^{-\lambda t}$ and
\[
\Pr(a+Y>t)=\begin{cases}1, & t<a,\\ e^{-\mu(t-a)}, & t\ge a.\end{cases}
\]
Hence
\[
\E[T]=\int_0^a e^{-\lambda t}\,dt+\int_a^\infty e^{-\lambda t}e^{-\mu(t-a)}\,dt
=\frac{1-e^{-\lambda a}}{\lambda}+e^{-\lambda a}\int_0^\infty e^{-(\lambda+\mu)s}\,ds
\quad(s=t-a),
\]
so
\[
\boxed{\E[T]=\frac{1-e^{-\lambda a}}{\lambda}+\frac{e^{-\lambda a}}{\lambda+\mu}.}
\]
(Interpretation: if task A has not finished by time $a$, which happens with
probability $e^{-\lambda a}$, then by memorylessness the remaining time is the
minimum of independent $\operatorname{Expo}(\lambda)$ and $\operatorname{Expo}(\mu)$,
i.e.\ $\operatorname{Expo}(\lambda+\mu)$ with mean $\frac1{\lambda+\mu}$.)

\textbf{(b)} Ties have probability $0$. Condition on whether $X<a$:
\begin{itemize}
\item If $X<a$, task A finishes before task B even starts:
$\Pr(X<a)=1-e^{-\lambda a}$.
\item If $X\ge a$ (probability $e^{-\lambda a}$), by the memoryless property
$X-a\mid X\ge a\sim\operatorname{Expo}(\lambda)$, independent of $Y$. A finishes
first iff $X-a<Y$, and for independent exponentials
$\Pr(X-a<Y\mid X\ge a)=\frac{\lambda}{\lambda+\mu}$.
\end{itemize}
By the law of total probability,
\[
\boxed{\Pr(\text{A finishes before B})=\Pr(X<a+Y)
=1-e^{-\lambda a}+e^{-\lambda a}\frac{\lambda}{\lambda+\mu}
=1-\frac{\mu}{\lambda+\mu}e^{-\lambda a}.}
\]


\end{homeworkProblem}
\newpage
\begin{homeworkProblem}{20}
Let $Z\sim N(0,1)$. Find the value of $a>0$ that maximizes
\[
\Pr(a<Z<2a),
\]
and express the maximum probability using the standard Normal
CDF $\Phi$.

\medskip
\solution


Let $g(a)=\Pr(a<Z<2a)=\Phi(2a)-\Phi(a)$ for $a>0$, and let
$\varphi(z)=\frac1{\sqrt{2\pi}}e^{-z^2/2}$ be the standard Normal PDF. Then
\[
g'(a)=2\varphi(2a)-\varphi(a)=\frac{1}{\sqrt{2\pi}}\left(2e^{-2a^2}-e^{-a^2/2}\right)
=\frac{e^{-a^2/2}}{\sqrt{2\pi}}\left(2e^{-3a^2/2}-1\right).
\]
The factor $2e^{-3a^2/2}-1$ is strictly decreasing in $a>0$, equals $1>0$ at
$a=0$, and tends to $-1$ as $a\to\infty$. It vanishes exactly when
\[
e^{-3a^2/2}=\frac12\iff \frac{3a^2}{2}=\ln 2\iff a^*=\sqrt{\frac{2\ln2}{3}}\approx0.6798.
\]
So $g'(a)>0$ on $(0,a^*)$ and $g'(a)<0$ on $(a^*,\infty)$: $g$ increases and then
decreases, so $a^*$ is the unique global maximizer on $a>0$ (also, $g(a)\to0$ as
$a\to0^+$ and as $a\to\infty$). The maximum probability is
\[
\boxed{\max_{a>0}\Pr(a<Z<2a)=\Phi\!\left(2\sqrt{\tfrac{2\ln2}{3}}\right)-\Phi\!\left(\sqrt{\tfrac{2\ln2}{3}}\right)
\approx\Phi(1.3596)-\Phi(0.6798)\approx 0.1613.}
\]

\end{homeworkProblem}
\newpage
\begin{homeworkProblem}{20}

Let $X_1,\ldots,X_{100}$ be independent and identically distributed
random variables with
\[
\Pr(X_i=-2)=\Pr(X_i=1)=\Pr(X_i=4)=\frac13.
\]
Use the Central Limit Theorem with an appropriate continuity
correction to approximate
\[
\Pr\left(90<\sum_{i=1}^{100}X_i<111\right).
\]
Explain how you determine the endpoints used in the continuity
correction. You may express your answer using the standard Normal
CDF $\Phi$.

\medskip
\solution


\textbf{Mean and variance.}
$\E[X_i]=\frac{-2+1+4}{3}=1$, $\E[X_i^2]=\frac{4+1+16}{3}=7$, so
$\Var(X_i)=7-1=6$. Let $S=\sum_{i=1}^{100}X_i$. Then
\[
\E[S]=100,\qquad \Var(S)=600,\qquad \sqrt{\Var(S)}=\sqrt{600}=10\sqrt6\approx24.495.
\]

\textbf{Lattice of possible values.} Each $X_i\in\{-2,1,4\}$ can be written as
$X_i=1+3W_i$ with $W_i\in\{-1,0,1\}$. Hence
\[
S=100+3\sum_{i=1}^{100}W_i,
\]
so $S$ only takes values in $\{\dots,88,91,94,\dots\}$, i.e.\ integers with
$S\equiv1\pmod 3$; neighbouring possible values are $3$ apart (not $1$).

\textbf{Endpoints of the continuity correction.} The possible values with
$90<S<111$ are exactly $S\in\{91,94,\dots,106,109\}$. The nearest possible
values outside the range are $88$ (below) and $112$ (above). The continuity
correction replaces each lattice point by an interval of width equal to the
lattice spacing $3$, centred at that point, so we cut halfway between the last
excluded and first included lattice points:
\[
\text{lower endpoint}=\frac{88+91}{2}=89.5,\qquad
\text{upper endpoint}=\frac{109+112}{2}=110.5.
\]
(Using $90.5$ and $110.5$, i.e.\ a correction of $\pm\frac12$, would be wrong here
because the spacing is $3$, not $1$; $90$ and $111$ themselves are not attainable.)
Thus
\[
\Pr(90<S<111)=\Pr(89.5<S<110.5).
\]

\textbf{CLT approximation.} By the CLT, $S\overset{\cdot}{\sim}N(100,600)$, so
\[
\Pr(90<S<111)\approx
\Phi\!\left(\frac{110.5-100}{\sqrt{600}}\right)-\Phi\!\left(\frac{89.5-100}{\sqrt{600}}\right)
=2\Phi\!\left(\frac{10.5}{10\sqrt6}\right)-1
\approx 2\Phi(0.4287)-1\approx 0.332.
\]

\end{homeworkProblem}
\end{document}